\section{Methodology}
\label{sec:methodology}

Building on our observation that repository UX may systematically shape documentation outcomes, we designed a measurement approach with four components: (1) friction scoring protocol for objective platform comparison, (2) large-scale automated extraction for statistical power, (3) binary presence detection for reproducibility, and (4) within-platform validation for causal evidence.

\subsection{Friction Score Definition}

We introduce \textbf{friction score} (0-4 scale) as a quantifiable UX metric capturing repository design quality through binary feature detection:

\begin{itemize}
\item \textbf{Automated extraction} (0/1): Platform provides automated field population from dataset analysis (e.g., OpenML's ARFF parser auto-extracts feature statistics)
\item \textbf{Pre-filled templates} (0/1): Platform provides structured templates with example content (e.g., HuggingFace dataset card sections pre-populated with guidance)
\item \textbf{Validation feedback} (0/1): Platform provides real-time validation with actionable error messages during metadata entry
\item \textbf{Programmatic API} (0/1): Platform supports programmatic dataset upload with metadata specification via API clients
\end{itemize}

The friction score is the binary sum of these features. This design ensures objectivity (features are documentable from platform APIs and documentation), reproducibility (binary detection eliminates subjective thresholds), and cross-platform comparability (platform-agnostic metric).

\textbf{Platform Friction Scores:}
\begin{itemize}
\item \textbf{UCI ML Repository} (friction=0): Manual web forms only, no automated extraction, no templates, no validation, no API
\item \textbf{OpenML} (friction=2): Automated ARFF extraction (1) + programmatic API via \texttt{openml-python} library (1)
\item \textbf{HuggingFace Datasets} (friction=3): Pre-filled dataset card templates (1) + validation feedback for required fields (1) + programmatic API via \texttt{datasets} library (1)
\end{itemize}

This 0-2-3 distribution provides variance for detecting friction-completeness correlation while controlling for confounds through stratified sampling.

\subsection{Large-Scale Metadata Extraction}

We extracted metadata from 10,000+ datasets across three platforms using platform-specific methods:

\textbf{OpenML (2,500 datasets):} API client via \texttt{openml-python} library, batch requests for dataset metadata, XML parsing for field extraction.

\textbf{HuggingFace (7,000 datasets):} API client via \texttt{huggingface\_hub.list\_datasets()}, YAML parsing of dataset card frontmatter, stratified random sampling proportional to platform size.

\textbf{UCI ML Repository (500 datasets):} Web scraping via BeautifulSoup, HTML table column parsing, smaller sample due to limited dataset count ($\sim$600 total).

\textbf{Temporal Snapshot:} All extraction completed within 2-week window (2026-08-19 snapshot) to eliminate temporal drift confounds. Single-timepoint design ensures apples-to-apples comparison without platform evolution interference.

\textbf{Stratified Sampling:} Proportional to platform size (OpenML $\sim$20k datasets, HuggingFace $\sim$60k datasets, UCI $\sim$600 datasets) ensures representativeness while maintaining statistical power. Power analysis: $n=10{,}000$, $\alpha=0.05$, effect size $d=0.5$ $\rightarrow$ power $>0.99$ for detecting 20pp+ differences.

\subsection{Binary Presence Detection}

We applied parsing rules for 6 target metadata fields:

\textbf{Optional Fields (not enforced):}
\begin{itemize}
\item \texttt{preprocessing\_code}: Present if $>$50 characters + language keywords (python, R, julia) OR file extension (.py, .R, .ipynb)
\item \texttt{data\_source\_url}: Present if valid URL pattern (http/https) excluding placeholder text
\item \texttt{collection\_date}: Present if date format (YYYY-MM-DD, YYYY-MM, YYYY) or temporal description
\end{itemize}

\textbf{Required Fields (platform-enforced):}
\begin{itemize}
\item \texttt{license}: Present if $>$5 characters excluding placeholders (e.g., ``None'', ``N/A'', ``Unknown'')
\item \texttt{version}: Present if semantic version (X.Y.Z), integer version (v1, v2), or auto-generated hash
\end{itemize}

\textbf{Cross-Platform Schema Normalization:} We documented explicit semantic mapping protocol: OpenML \texttt{original\_data\_url} $\rightarrow$ HuggingFace \texttt{source\_url} $\rightarrow$ UCI \texttt{data source link} normalized to \texttt{data\_source\_url}. Finite platform set (3) makes manual mapping tractable (validated via 50-sample manual inspection showing 82.7\% semantic agreement).

\textbf{Validation:} 100-sample stratified manual review confirmed parsing accuracy $>$90\%, meeting feasibility threshold. Binary detection eliminates subjective quality thresholds while maintaining reproducibility.

\subsection{Within-Platform Causal Validation}

To control for platform-level confounds (age, funding, community size), we conducted within-HuggingFace comparison:

\textbf{Upload Method Proxy:} API-uploaded datasets identified via programmatic metadata patterns (auto-generated README structure, API-specific field formats, commit history analysis). Manual-uploaded datasets identified via web form patterns (human-written README, form-specific metadata structure).

\textbf{Sample:} 400 datasets (200 API, 200 manual) stratified random sampling from HuggingFace corpus.

\textbf{Metric:} Metadata completeness score (0-100\%) calculated as mean presence across 6 target fields.

\textbf{Statistical Test:} Welch $t$-test (two-sided) for mean completeness difference, normality checks via Shapiro-Wilk, Cohen's $d$ effect size.

This within-platform design provides stronger causal evidence than cross-platform comparison alone, as it eliminates confounds from platform age, funding model, and community demographics.

\subsection{Experimental Design Summary}

\begin{table}[h]
\centering
\caption{Methodology design choices and rationale}
\label{tab:methodology_design}
\begin{tabular}{lll}
\toprule
Component & Choice & Rationale \\
\midrule
Friction score & 0-4 binary sum & Objective, documentable \\
Sample size & 10,000+ datasets & Power $>0.99$ for 20pp+ diffs \\
Presence detection & Binary parsing & Automated, reproducible \\
Within-platform & HF API vs manual & Controls confounds \\
Temporal design & Single snapshot & Eliminates drift \\
\bottomrule
\end{tabular}
\end{table}

Our methodology addresses prior work limitations: Yang (2024) analyzed single platform, we compare 3 platforms; Strecker (2026) used qualitative taxonomy, we measure quantitative presence rates; prior work lacked explicit friction measurement, we introduce friction score as operationalization of UX quality.
